% Generated draft: requires mathematical and source review.
\documentclass[11pt]{article}
\usepackage[margin=25mm]{geometry}
\usepackage{fontspec}
\setmainfont{Latin Modern Roman}
\usepackage{amsmath,amssymb,mathtools,unicode-math}
\setmathfont{Latin Modern Math}
\usepackage{enumitem,needspace,xurl,hyperref}
\setlength{\parindent}{0pt}
\setlength{\parskip}{5pt}
\setlength{\emergencystretch}{3em}
\newcommand{\N}{\mathbb{N}}
\newcommand{\Z}{\mathbb{Z}}
\newcommand{\Q}{\mathbb{Q}}
\newcommand{\R}{\mathbb{R}}
\newcommand{\C}{\mathbb{C}}
\newcommand{\F}{\mathbb{F}}
\newcommand{\A}{\mathbb{A}}
\newcommand{\PP}{\mathbb{P}}
\newcommand{\E}{\mathbb{E}}
\newcommand{\eps}{\varepsilon}
\newcommand{\dd}{\,\mathrm{d}}
\begin{document}
\section*{UnsolvedMath ID: 6000009}
\textbf{Source problem code:} catalogue entry
\subsection{Definitions and mathematical statement}
A statistical manifold in the usual Riemannian convention is \((M,h,\nabla)\), where \(h\) is a Riemannian metric, \(\nabla\) is torsion-free, and \(\nabla h\) is totally symmetric. Fix the curvature convention
\[
R(X,Y)Z=\nabla_X\nabla_YZ-\nabla_Y\nabla_XZ-\nabla_{[X,Y]}Z,
\qquad K(X,Y,Z,W)=h(R(X,Y)Z,W).
\]
At a point let \(V\) be the real tangent vector space. The algebraic tensor space specified by Kurose's question is
\[
\mathcal K_{\mathrm{stat}}(V)=\{K\in(V^*)^{\otimes4}:\text{ the following identities hold for all }X,Y,Z,W\in V\}:
\]
\[
K(X,Y,Z,W)=-K(Y,X,Z,W),
\]
\[
K(X,Y,Z,W)+K(Y,Z,X,W)+K(Z,X,Y,W)=0,
\]
\[
K(X,Y,Z,W)+K(Y,W,Z,X)+K(W,X,Z,Y)=0.
\]
The middle equation is the first Bianchi identity. The last equation retains exactly the source's argument positions. No extra last-pair antisymmetry or interchange symmetry is imposed.

Investigate the structure of this linear space in each dimension. The identities are basis-independent, so it is naturally a representation under change of basis in \(V\); with a fixed metric one can instead ask for its orthogonal-group decomposition and contractions. The source does not choose a particular output such as a dimension formula, an irreducible decomposition under a specified group, or a geometric realization theorem. In particular, classification of the algebraic tensors is not silently replaced by classification of statistical manifolds realizing them.

The supplied research summary contrasts the standard metric algebraic-curvature decomposition with the less symmetric statistical setting and labels the general structural problem partially understood. It provides no explicit decomposition or verified definitive classification to add to the stated identities.
\subsection{Short English statement}
Describe the space of four-covariant tensors satisfying the three displayed curvature identities, without adding the symmetries of Riemannian curvature. The intended type of structural classification remains to be specified.
\subsection{Sources}
\begin{itemize}
\item Local UnsolvedMath registry, permanent ID \(6000009\), exact source code \texttt{catalogue entry}, ``Curvature-Type Tensors of Statistical Manifolds.'' Proposer Takashi Kurose, 1998.
\item Furuhata, Matsuzoe and Urakawa, \emph{Open Problems in Affine Differential Geometry and Related Topics} (1998), item 3(c), recorded scanned PDF pages 1--2; \url{https://www.jstage.jst.go.jp/article/iis/4/2/4_2_125/_pdf/-char/en}. Recorded access 2026-07-29; manual OCR from page images.
\item Supplied PARTIAL-PROGRESS research summary discusses metric versus statistical curvature decomposition without giving a further precise publication citation or decomposition theorem.
\end{itemize}
\end{document}
